\chapter*{Nota de actualización, 9 de septiembre de 2026}
\addcontentsline{toc}{chapter}{Nota de actualización, 9 de septiembre}

\titular{Los anexos de la cola del paquete \textbf{no se sirvieron hasta el día siguiente}
de la entrega, y el índice de la Gaceta \textbf{reasignó las letras} entretanto. Con ellos
llega una vara que este documento no tenía: la Secretaría de Bienestar estima en
\textbf{226.133,3 mdp} ---\textbf{0,57~\% del PIB}--- lo necesario para cumplimentar la
política de subsidios de vivienda y suelo en 2027. El \textbf{Ramo 15 completo} ---que es
mucho más que subsidios de vivienda--- son \textbf{37.873,6 mdp}: el \textbf{16,7~\%} de esa
cifra.}

\textbf{Qué es esta nota y qué no es.} El cuerpo de este documento se escribió el 8 de
septiembre de 2026, entre las 18:45 y las 21:00, con el paquete que existía esa tarde. No se
ha reescrito. Esta nota se añade el día siguiente y \textbf{no modifica ninguna cifra ni
ningún capítulo}: registra lo que llegó después, corrige un renglón del registro que quedó
mal contado, y publica dos hallazgos con cifra que el paquete completo permite y el paquete
del día 8 no permitía. \textbf{Lo que estaba fuera de alcance sigue fuera de alcance.}

\section*{Uno. Lo que no cambió: la Ruta B sigue vigente}

Se volvieron a sondear los seis recursos que bloquearon la corrida. \textbf{Ninguno cambió
de estado} entre las 21:45 del día 8 y las 07:33 del día 9: los cuatro analíticos del
proyecto siguen en \textbf{404} y las dos ranuras del portal de datos abiertos en
\textbf{503}. La ranura de plazas sigue sirviendo, sin aviso, el archivo del ejercicio 2026.

En consecuencia, \textbf{todo lo que los dieciséis capítulos declaran como no afirmable lo
sigue siendo}: no hay diff institucional, no hay auditoría de etiquetado de los anexos
transversales, no hay composición por programa presupuestario, no hay distribución
territorial, y la separación entre atención médica y pensiones en el IMSS y el ISSSTE sigue
sin poder hacerse.

\section*{Dos. Las letras de la Gaceta se reasignaron después de la entrega}

El registro del documento anota que los anexos F y G ---que el índice del 8 de septiembre
describía como la iniciativa de Código Fiscal y el informe arancelario--- estaban enlazados
y respondían 404. \textbf{El 404 era cierto; la descripción no.}

Al reintentarlos el día 9 responden, pero con otro contenido. El índice de la Gaceta se
reescribió: \textbf{la iniciativa de Código Fiscal desapareció del índice sin haberse
servido nunca}, la letra F pasó a ser una reforma a la \textbf{Ley Aduanera}, se insertaron
dos iniciativas que no estaban ---la \textbf{Ley de Economía Digital para Pagos Digitales y
Electrónicos} y la \textbf{Ley General para el Fortalecimiento y Armonización Catastral y
Registral}--- y \textbf{el informe arancelario del artículo 131 se recorrió de la G a la J}.
Además se sirven cuatro anexos más, K a N, que \textbf{el índice todavía no lista}: sólo
aparecen en un bloque comentado dentro del HTML, con las letras del año anterior y fechados
2026, cuando los PDF dicen 2027.

\begin{table}[htbp]\centering\footnotesize
\caption{Los trece anexos del paquete que sirve la Gaceta 7121, y la letra que el índice les daba el día de la entrega}\label{cua:anexos}
\begin{tabular}{l>{\raggedright\arraybackslash}p{5.0cm}rl>{\raggedright\arraybackslash}p{2.9cm}}
\toprule
Letra & Contenido, leído de la primera página del PDF & pp. & Creado & Letra en el índice del 8 sep. \\
\midrule
A & ILIF 2027 & 310 & 8 sep. & A \\
B & PPEF, proyecto de decreto & 254 & 8 sep. & B \\
C & CGPE 2027 & 78 & 8 sep. & C \\
D & Ley Federal de Derechos & 142 & 8 sep. & D \\
E & Ley del Impuesto sobre la Renta & 228 & 8 sep. & E \\
\textbf{F} & Ley Aduanera & 28 & 8 sep. 19:18 & no listada; la F anunciaba Código Fiscal \\
\textbf{G} & Ley de Economía Digital para Pagos Digitales y Electrónicos & 22 & 8 sep. 19:06 & no listada \\
\textbf{H} & Ley General para el Fortalecimiento y Armonización Catastral y Registral & 28 & 8 sep. 19:08 & no listada \\
\textbf{J} & Informe arancelario, artículo 131 & 12 & 8 sep. 19:54 & G \\
\textbf{K} & Nota metodológica y declaratoria ZAP 2027 & 18 & 9 sep. 05:29 & no listada \\
\textbf{L} & Listado ZAP rurales 2027 & 244 & 9 sep. 05:38 & no listada \\
\textbf{M} & Listado ZAP urbanas 2027 & 224 & 9 sep. 05:41 & no listada \\
\textbf{N} & Estimación de subsidios de vivienda y suelo 2027 & 6 & 9 sep. 05:46 & no listada \\
\bottomrule
\end{tabular}
\\[2pt]\begin{minipage}{0.92\textwidth}\scriptsize Leído el 9 de septiembre de 2026 a las 07:40; el paginado es el de la edición de la Gaceta, que añade portada y colofón. \textbf{La letra no identifica el contenido}: el índice del día de la entrega se reescribió y las letras se reasignaron. Los anexos K a N no aparecen en el índice vivo. Las dos copias del índice, del 8 y del 9, están archivadas.\end{minipage}
\end{table}


Los anexos A, C, D y E se volvieron a descargar y son \textbf{idénticos byte a byte} a los
que sostienen este documento. \textbf{Lo que se movió es la cola del paquete, no la cabeza:
ninguna cifra de los dieciséis capítulos depende de un anexo reasignado.}

\textbf{Una corrección de método, que vale más que el episodio.} La advertencia con la que
este proyecto llegó a la entrega era «las letras no son fijas: léelas del índice». No basta.
\textbf{El índice mismo se reescribe, y las letras se reasignan.} La regla que sustituye a la
anterior: la letra no identifica nada; identifica el PDF por su primera página, guarda el
índice con marca de tiempo cada vez que lo consultes, y recorre las letras de la cola hasta
el primer 404 firme. Se anota también que \textbf{el 404 de la Gaceta es intermitente}: el
anexo L respondió 404 a las 07:34 y 206 a las 07:38, misma dirección, cuatro minutos después.
El almacén no retira archivos: los sirve tarde y de forma inestable.

\section*{Tres. Lo que llegó, y qué de ello es utilizable}

\textbf{Las tres iniciativas de ley y el informe arancelario vienen escaneados.} Los anexos
F, G, H y J tienen capa de texto \textbf{sólo en la portada y el colofón}: 2 de 28 páginas,
2 de 22, 2 de 28 y 2 de 12. El cuerpo es imagen. \textbf{Se archivan y no se explotan aquí}:
el reconocimiento óptico de un texto legal no se usa para citar cifras sin cotejo a mano, y
este documento no cita lo que no ha leído.

Los cuatro anexos de la Secretaría de Bienestar sí traen texto completo, y de ellos salen los
dos hallazgos que siguen.

\subsection*{3.1 La vara de suficiencia en vivienda}

El artículo 61 de la Ley de Vivienda obliga a la Secretaría de Bienestar a estimar cada año
el monto de recursos federales requeridos para cumplimentar la política de subsidios en
materia de vivienda y suelo. \textbf{Es una vara externa al presupuesto, publicada en el
mismo paquete}, y es exactamente el tipo de contraste que este proyecto busca y que en salud
lleva tres ejercicios sin poder hacer.

\begin{table}[htbp]\centering\footnotesize
\caption{La vara del artículo 61 de la Ley de Vivienda contra el Ramo 15 (mdp)}\label{cua:vivienda}
\begin{tabular}{lrrrr}
\toprule
Necesidad & Hogares & UMA & Costo unitario (pesos) & Monto (mdp) \\
\midrule
Mejora de vivienda & 849.226 & 25 & 89.155,5 & 75.713,2 \\
Ampliación de vivienda & 771.091 & 50 & 178.311,0 & 137.494,0 \\
Vivienda nueva & 36.246 & 100 & 356.622,0 & 12.926,1 \\
\midrule \textbf{Requerimiento declarado} & & & & \textbf{226.133,3} \\
Ramo 15 completo, proyecto 2027 & & & & 37.873,6 \\
\textbf{El ramo como proporción del requerimiento} & & & & \textbf{16,7\,\%} \\
\bottomrule
\end{tabular}
\\[2pt]\begin{minipage}{0.92\textwidth}\scriptsize Fuente: Gaceta Parlamentaria 7121, anexo N, estimación de la Secretaría de Bienestar; y \texttt{datos/ramos.csv} para el Ramo 15. \textbf{Años de pesos distintos}: los costos unitarios provienen de las Reglas de Operación y de la UMA de 2026, y el presupuesto está en pesos de 2027. \textbf{El Ramo 15 no es el perímetro del gasto en subsidios de vivienda}: es el ramo completo, y CONAVI es un organismo descentralizado. La comparación es de orden de magnitud.\end{minipage}
\end{table}


El requerimiento declarado es de \textbf{226.133,3 mdp}: 849.226 hogares que requieren
mejora de vivienda, 771.091 que requieren ampliación y 36.246 que requieren vivienda nueva,
costeados a 25, 50 y 100 veces la UMA mensual. Contra eso, el \textbf{Ramo 15, Desarrollo
Agrario, Territorial y Urbano, completo}, son \textbf{37.873,6 mdp} en el proyecto 2027, con
un crecimiento real de \textbf{0,35~\%} respecto del proyecto 2026.

\porqueimporta{\textbf{El requerimiento es seis veces el ramo entero}, y el ramo entero es
mucho más que subsidios de vivienda. Dicho en la unidad que este documento usa para todo lo
demás: la política declarada necesaria vale \textbf{0,57~\% del PIB} y el ramo entero que
la aloja, junto con todo lo demás que ese ramo hace, vale \textbf{0,10~\%}. \textbf{La brecha no es un
matiz de calibración: es de orden de magnitud}, y está declarada por la misma dependencia,
en el mismo paquete, el mismo día.}

\textbf{Tres advertencias, sin las cuales la cifra no debe citarse.} Primera: los costos
unitarios salen de las Reglas de Operación \textbf{de 2026} y de la \textbf{UMA de 2026}, de
modo que el requerimiento está en pesos de 2026 y el presupuesto en pesos de 2027; corregirlo
ampliaría la brecha, y este documento \textbf{no lo corrige} porque la fuente no publica la
actualización. Segunda: \textbf{CONAVI es un organismo descentralizado} y no todo el gasto en
vivienda pasa por el Ramo 15; sin el analítico por programa ---que sigue en 404--- no puede
construirse el perímetro exacto del gasto en subsidios de vivienda. Tercera, en consecuencia:
\textbf{esto es una comparación de orden de magnitud a nivel de ramo, no un cierre contable},
y así debe leerse.

\subsection*{3.2 La declaratoria de zonas de atención prioritaria 2027}

Las ZAP rurales de 2027 quedan integradas por \textbf{1.575 municipios} en las 32 entidades
federativas; las urbanas, por \textbf{43.636 áreas geoestadísticas básicas} distribuidas en
\textbf{4.531 localidades urbanas de 2.423 municipios}. La nota metodológica documenta la
construcción criterio por criterio y de forma acumulativa: 790 municipios por grado de
marginación, 796 al añadir rezago social, 953 al añadir municipios indígenas, 959 al añadir
población afromexicana, 1.383 al añadir incidencia delictiva, 1.621 al añadir pobreza
extrema, 1.638 al añadir grado de accesibilidad a carretera pavimentada y 1.645 al incorporar
municipios de creación reciente, de los que 70 pasan a la lista urbana por ser urbanos.

\porqueimporta{\textbf{Es el perímetro territorial que se declara cada año junto con el
presupuesto, publicado con su regla de construcción a la vista.} Este documento no lo cruza con el gasto
---hacerlo exige la distribución territorial, que la Ruta B no permite--- pero lo deja
registrado como denominador disponible para el ejercicio siguiente. \textbf{La nota
metodológica publica además una ruta de descarga de las 43.636 AGEB con sus variables}, que
la precarga de este proyecto no tenía.}

\section*{Cuatro. Qué se vuelve a correr y qué no}

Esta nota \textbf{no reescribe ningún capítulo}, de modo que las verificaciones que validan
el cuerpo siguen siendo válidas tal como se publicaron. Las cifras que la nota introduce
---las de los cuadros de esta sección--- se verificaron contra su fuente y se incorporaron al
guion de identidades, que pasa de \textbf{29 a 41 pruebas y sigue cerrando sin una sola
falla}: los tres montos de vivienda contra hogares por costo unitario, los tres desgloses
urbano-rural, la UMA implícita en los tres costos, las dos sumas del total ---en pesos y en
UMA--- y la aritmética de los municipios de la ZAP. La checklist de omisiones incorpora dos
renglones nuevos, y la vara de vivienda queda anotada como \textbf{resuelta}. La compuerta
semántica se aplicó al titular, a los dos pies de cuadro y a las dos frases de «por qué
importa» de esta nota: \textbf{tres frases reescritas, ninguna cifra modificada}, todas por
sujeto más ancho que su cálculo. \textbf{Ninguna cifra del cuerpo del documento cambió.}

\textbf{Un error de conteo del cuerpo, corregido.} El registro decía que la compuerta se
aplicó a los pies de «los diecinueve cuadros». Los cuadros del documento del día 8 eran
\textbf{veinte}, y con los dos de esta nota son veintidós. Se corrigió el número; no cambia
ninguna cifra ni ninguna afirmación.

\begin{ficha}{nota de actualización}
\fila{Perímetro}{Ramo 15 completo, del Anexo 1 del decreto. \textbf{No es el perímetro del
gasto en subsidios de vivienda}, que no puede construirse sin el analítico por programa.}
\fila{Línea de comparación}{Línea P, proyecto 2026 contra proyecto 2027, para el Ramo 15.
Para el requerimiento no hay línea: es el primer año que este proyecto lo recoge.}
\fila{Deflactor}{Del PIB, 1,040, del CGPE 2027, en la variación real del Ramo 15.}
\fila{Año de los pesos}{\textbf{Mixto, y es el defecto de la comparación}: el requerimiento
está en pesos de 2026 (UMA y Reglas de Operación de 2026) y el presupuesto en pesos de 2027.}
\fila{PIB y añada}{PIB nominal 2027 de 39.419.400,0 mdp, CGPE 2027 edición SHCP, Anexo II.5.}
\fila{Denominadores}{Ninguno per cápita. Los conteos de hogares y de municipios se citan como
los publica la fuente.}
\fila{Fuentes}{Gaceta Parlamentaria 7121, anexos K (nota metodológica y declaratoria ZAP
2027), L y M (listados rural y urbano) y N (estimación del artículo 61 de la Ley de
Vivienda); \texttt{datos/ramos.csv} para el Ramo 15; las dos copias del índice de la Gaceta,
del 8 y del 9 de septiembre, para la reasignación de letras.}
\fila{Qué no puede afirmarse con ellas}{Que el presupuesto de vivienda cubra el 16,7~\% de la
necesidad. Lo que se compara es \emph{un ramo completo} contra \emph{un requerimiento de
política}, en años de pesos distintos y con parte del gasto fuera del ramo. La afirmación
sostenible es de orden de magnitud: \textbf{el requerimiento declarado es varias veces el
ramo entero}.}
\end{ficha}
